% ============================================================
% Formulario, argomento 3: la probabilità e le sue regole (dispensa, capitoli 3 e 4)
% ============================================================
\argomento{Probabilità: definizioni e regole}{Dispensa: capitoli 3 e 4}

\voce{frequenza}{Frequenza relativa}{$n_A$ = volte in cui si verifica $A$ su $n$ prove}
  {f_n(A) = \frac{n_A}{n} \qquad 0 \leq f_n(A) \leq 1}
  {}

\voce{frequentista}{Probabilità come limite della frequenza}{definizione frequentista}
  {P(A) = \lim_{n \to \infty} \frac{n_A}{n}}
  {}

\voce{laplace}{Spazi equiprobabili}{tutti i risultati alla pari}
  {P(A) = \frac{\card{A}}{\card{\Omega}} = \frac{\text{casi favorevoli}}{\text{casi possibili}}}
  {\variante{un solo risultato}{P(\{\omega\}) = \frac{1}{\card{\Omega}}}}

\voce{assiomi}{Assiomi di Kolmogorov}{legge di probabilità $P : \mathcal{E} \to [0, 1]$}
  {P(A) \geq 0 \qquad P(\Omega) = 1 \qquad A \cap B = \emptyset \;\Rightarrow\; P(A \cup B) = P(A) + P(B)}
  {\variante{numerabile additività, $B_i \cap B_j = \emptyset$ per $i \neq j$}
            {P\left( \bigcup_{n=1}^{\infty} B_n \right) = \sum_{n=1}^{\infty} P(B_n)}}

\voce{pesi}{Probabilità su uno spazio discreto}{un peso $p(\omega)$ per ogni risultato}
  {P(A) = \sum_{\omega \in A} p(\omega) \qquad p(\omega) \geq 0, \quad \sum_{\omega \in \Omega} p(\omega) = 1}
  {}

\voce{compl}{Evento complementare}{}
  {P(\compl{A}) = 1 - P(A)}
  {\variante{evento impossibile}{P(\emptyset) = 0}
   \variante{«almeno uno»}{P(\text{almeno uno}) = 1 - P(\text{nessuno})}}

\voce{unione}{Unione di due eventi}{}
  {P(A \cup B) = P(A) + P(B) - P(A \cap B)}
  {\variante{$A$ e $B$ incompatibili}{P(A \cup B) = P(A) + P(B)}
   \variante{sub-additività}{P(A \cup B) \leq P(A) + P(B)}
   \variante{$M$ eventi a due a due incompatibili}{P\left( \bigcup_{i=1}^{M} A_i \right) = \sum_{i=1}^{M} P(A_i)}
   \variante{tre eventi \oltre}{{P(A \cup B \cup C)} = {P(A)} + {P(B)} + {P(C)} - {P(A \cap B)} - {P(A \cap C)}
                                 - {P(B \cap C)} + {P(A \cap B \cap C)}}}

\voce{sottrazione}{Differenza (sottrazione) fra eventi}{probabilità di «$A$ ma non $B$»}
  {P(A \setminus B) = P(A \cap \compl{B}) = P(A) - P(A \cap B)}
  {\variante{se $B \subseteq A$}{P(A \setminus B) = P(A) - P(B)}}

\voce{monotonia}{Monotonia e limiti della probabilità}{}
  {A \subseteq B \;\Rightarrow\; P(A) \leq P(B) \qquad\qquad 0 \leq P(A) \leq 1}
  {\variante{evento certo e impossibile}{P(\Omega) = 1, \quad P(\emptyset) = 0}}
